% ===========================================================================
% SECTION 7 — DISCUSSION
% Target: 600-800 words
% ===========================================================================

\section{Discussion}
\label{sec:discussion}

\subsection{Main Findings}
\label{subsec:main_findings}

Four findings emerge from this study. First, a systematic six-architecture
comparison (Section~\ref{sec:ml_surrogates}) shows that prediction accuracy
alone is a poor predictor of closed-loop usefulness: LLNFM and the two GP
variants dominate on one-step RMSE, yet PINN — a mid-table predictor —
achieves the best V1 closed-loop RMSE (Section~\ref{subsec:v1_results}),
while GP-v2's substantially improved open-loop accuracy over GP V1 leaves
closed-loop RMSE unchanged (Section~\ref{subsec:v1_v2_comparison}). Second,
what initially appeared to be a hard operational incompatibility between
LSTM, TCN, SW-MLP, PINN-v2, and GP-v2 on one hand and closed-loop
SQP-based \texttt{nlmpc} on the other, was traced through a five-stage
diagnostic (Section~\ref{subsec:operational_feasibility}) to a specific,
resolvable cause: repeated calls to MATLAB's \texttt{predict()} function
inside the SQP solve, not surrogate architecture, model size, or MATLAB
version. A verified manual forward pass --- explicit matrix arithmetic
for feedforward and Mat\'ern-kernel architectures, and a vectorized
causal convolution for TCN --- restored closed-loop viability for all
six surrogates, with speedups of $7.8\times$ to $2844\times$ and
near-elimination of real-time budget overruns
(Table~\ref{tab:predict_bypass}). This reframes what the ML-MPC-for-wind
literature would otherwise record as an architectural limitation of
deep-learning surrogates as, instead, an implementation-specific cost
that is both diagnosable and remediable within the same toolchain,
consistent with \citet{salzmann_real-time_2022}'s independent report
that decoupling a neural network's evaluation from the surrounding
solver's iteration loop is necessary --- though not previously
documented, to our knowledge, as arising specifically from
\texttt{predict()}'s own per-call overhead in MATLAB's \texttt{nlmpc}.
A structurally related, but functionally distinct, finding is reported by
\citet{witter_creating_2025}: their LSTM/GRU surrogates diverge in closed
loop for a \emph{predictive-accuracy} reason (the learned dynamics
themselves behave differently once embedded in a feedback loop), whereas
our LSTM's closed-loop difficulty was, before the bypass,
\emph{computational} (excessive per-step latency causing the solver to
operate on a stale or default control action) rather than a divergence of
the learned dynamics --- the extended Monte Carlo evaluation
(Section~\ref{subsec:monte_carlo}) shows LSTM achieving competitive, and
at some wind speeds the best, tracking accuracy once the computational
bottleneck is removed. Distinguishing these two failure modes ---
functional divergence versus computational infeasibility --- is
important for practitioners, since only the latter admits the kind of
implementation-level remedy demonstrated here. Third,
the Region~II/III torque switching correction
(Section~\ref{subsec:torque_switching}) — while standard practice in the
classical wind turbine control literature \citep{bianchi2006,pao2009} — was
in our case discovered only through long-horizon closed-loop simulation
with an incorrect constant-torque assumption; the fact that a textbook
modeling choice reproduces a documented irrecoverable failure mode
underscores that this detail cannot be treated as a minor implementation
nicety when generating training data or evaluating closed-loop MPC.
Fourth, calibrating GP hyperparameters by marginal-likelihood optimization
under a Matérn~5/2 kernel \citep{rasmussen2006} yields a predictive
uncertainty that is usable directly as an MPC cost penalty
(Section~\ref{subsec:gp_uncertainty_cost}), in the spirit of cautious /
uncertainty-aware MPC \citep{hewing2020,kocijan2004} but applied here, to
our knowledge, for the first time to a Region~III wind turbine speed
regulation problem with an ML surrogate as the internal predictor.
Relative to the model-free reinforcement-learning alternative to
surrogate-based MPC \citep{xie_data-driven_2023}, our approach retains
MPC's explicit constraint handling and, via the manual bypass of
Section~\ref{subsec:operational_feasibility}, its real-time feasibility
argument grounded in measured per-step cost rather than only
asymptotic policy performance; a direct empirical comparison between
the two paradigms on the same plant is left to future work.

\subsection{Limitations}
\label{subsec:limitations}

Several limitations qualify these findings. The plant model retains only
two states ($\omegar,\beta$) and omits drivetrain torsion and tower
fore-aft dynamics, both of which are standard in higher-fidelity wind
turbine control models \citep{schlipf2013} and materially affect load-
relevant behavior; the reduced model is appropriate for a first
systematic surrogate comparison but likely understates the wind
turbine's dynamic order.
The wind input is spatially uniform (Section~\ref{subsec:wind_model}):
no rotational sampling or tower shadow effect is modeled, both of which
introduce periodic loading components a real turbine controller must
reject. The plant is integrated with explicit Euler at
$\Delta t=\SI{0.1}{s}$, equal to the pitch actuator time constant
$\tau_\beta$. A dedicated sensitivity check, run post hoc, confirms
this is a genuine limitation rather than the merely hypothetical
concern noted in Section~\ref{subsubsec:pinn_v1} for PINN's
degenerate residual: because $\Delta t=\tau_\beta$, Euler's discrete
pitch-actuator update reduces to $\beta_{k+1}=\beta_k+(u-\beta_k)=u$,
an instantaneous jump to the commanded pitch rather than the
$(1-e^{-1})\approx63\%$ single-time-constant approach a continuous
first-order lag would produce. Repeating a $\SI{60}{s}$
closed-loop run with the Baseline controller under classical RK4
integration of the identical continuous dynamics, instead of Euler,
leaves closed-loop RMSE essentially unchanged ($1.2709$ vs.\
$\SI{1.2734}{rpm}$, a $0.19\%$ relative difference), but cumulative
pitch activity drops by $54\%$ ($372.5$ vs.\ $\SI{172.0}{deg}$); a
follow-up variant isolating whether this arises from RK4's pitch-rate
saturation being evaluated once per integration substage (rather than
once per step, as in Euler) found the discrepancy essentially
unchanged when saturation semantics were matched to Euler's
($\SI{171.6}{deg}$), ruling out saturation handling as the cause and
confirming the Euler/$\tau_\beta$ degeneracy itself as responsible.
Because the same plant integrator underlies every controller and
surrogate reported in this paper, this bias is common across all
reported results rather than differentially favoring any one
architecture, so the \emph{relative} comparisons between controllers
that this paper's contributions rest on --- computational speedup,
relative tracking accuracy, and closed-loop feasibility verdicts ---
are not expected to be affected. However, cumulative pitch activity
figures reported throughout this paper (Table~\ref{tab:v1_results} and
subsequent tables) should accordingly be read as a comparative,
architecture-relative indicator rather than a physically calibrated
estimate of actuator wear; a smaller integration step decoupled from
$\tau_\beta$, noted as future work below, would be required before
absolute pitch-activity figures could be trusted on their own.
Whether the same integration artifact still explains the differing
closed-loop behavior of window-based surrogates once evaluated via the
manual bypass (Section~\ref{subsec:bypass_closed_loop}) is a related
question the present results do not settle. The
extended Monte Carlo evaluation (Section~\ref{subsec:monte_carlo}) covers
three wind speeds and five seeds each, but a single $\SI{60}{s}$ window
and no multi-objective load metric beyond rotor-speed RMSE and
cumulative pitch activity; it also surfaced an unresolved instability
(LSTM's coefficient of variation reaching $36.7\%$ at
$V=\SI{12}{m/s}$, versus $7$--$8\%$ at higher wind speeds) that this
study does not explain and that would need to be understood before
LSTM's otherwise competitive mean accuracy is trusted near the
Region~II/III boundary. The manual-bypass technique itself
(Section~\ref{subsec:operational_feasibility}) is a further limitation
in practice, not only a solution: each architecture required its own
hand-verified forward-pass implementation (layer-by-layer for
feedforward and convolutional networks, kernel-based for GP, gate
equations for LSTM), and the LSTM case specifically showed a
verification tolerance an order of magnitude looser
($10^{-4}$ vs.\ $10^{-7}$--$10^{-13}$) than the other five
architectures, attributed to single- versus double-precision
accumulation over recurrent steps rather than a logical error but not
independently confirmed; this per-architecture verification burden does
not disappear simply because the technique is shown to work here.
The V1-to-V2 evaluation discrepancy noted for nominally unchanged
components (Persistence, LLNFM) is, as established in
Section~\ref{subsec:comparative_accuracy}, fully traced to a
deliberate widening of the V2 dataset's initial-condition sampling to
cover Region~II, not an unresolved artifact. A genuinely open
methodological gap remains, however, around the GP-uncertainty cost
weight $\kappa$ (Section~\ref{subsec:gp_uncertainty_cost}): an attempt
to sweep $\kappa \in \{0, 0.2, 0.5, 0.8, 1.0\}$ under the original
\texttt{predict()}-based GP-v2 cost on a shortened $\SI{20}{s}$ window
returned a numerically identical closed-loop RMSE and zero pitch
activity for every $\kappa$ value tested, traced to
\texttt{ExitFlag}$<0$ (solver infeasibility) at every one of $200$
steps regardless of $\kappa$ — the pitch command never left its
initial value, so the flat RMSE across $\kappa$ reflects a shared
solver failure rather than genuine insensitivity to the uncertainty
penalty. This echoes, in a different form, the systematic
non-convergence under \texttt{MaxIterations}$=30$ documented for
GP-v2 in the Jacobian test of Section~\ref{subsec:operational_feasibility}
(Stage 3), and suggests the uncertainty-penalized cost may be
structurally harder for the SQP solver to satisfy than the standard
quadratic cost, independent of $\kappa$'s specific value or the
shortened horizon used for this attempt. Whether $\kappa=0.8$ is a
robust operating point therefore remains an open question, not a
settled one, and is carried forward as future work below rather than
claimed as validated here. Finally, prediction and tracking accuracy
are reported throughout this paper as RMSE alone; MAE and $R^2$ were
not computed retroactively for already-reported results, since doing
so would require the original per-sample prediction arrays rather than
the aggregate RMSE values actually logged during each run --- a
utility function (\texttt{compute\_extended\_metrics.m}) is provided
for future evaluation runs to report all three metrics jointly,
avoiding a repeat of this retroactive gap.

\subsection{Perspectives (V3)}
\label{subsec:perspectives_v3}

Five directions follow from the findings above. First, the LSTM
instability at $V=\SI{12}{m/s}$ (Section~\ref{subsec:monte_carlo}) and
the crossover in wind-speed sensitivity between
Baseline/GP/residual-MLP/SW-MLP and PINN-v2/TCN
(Section~\ref{subsec:monte_carlo}) both warrant targeted investigation
--- extending the wind-speed sweep beyond three points and pairing it
with the multi-objective load metrics noted above would clarify whether
these are genuine architectural properties or artifacts of the specific
training data distribution. Second, extending the plant to four states
by adding drivetrain torsion and tower fore-aft dynamics
\citep{schlipf2013} would let the comparison generalize to
load-relevant metrics, not only rotor-speed tracking, and would provide
a further, independent test of whether the manual-bypass technique
scales to a higher-order plant and correspondingly larger surrogates.
Third, decoupling the plant integration step from the pitch actuator
time constant ($\Delta t \ne \tau_\beta$, e.g.\ via RK4 integration or
a smaller $\Delta t$ with actuator sub-stepping) would remove the
Euler/$\tau_\beta$ degeneracy identified above and allow the cumulative
pitch-activity figures reported throughout this paper to be
re-evaluated as physically calibrated actuator-wear estimates, rather
than only architecture-relative comparisons. Fourth, an explicit
wind-speed forecasting or estimation layer — an
autoregressive predictor over the MPC horizon, or an Extended Kalman
Filter estimating rotor-effective wind speed from standard turbine
measurements \citep{soltani2013} — would let the controller anticipate
disturbances rather than react to the current measured $V$ alone, and
is a natural complement to the GP uncertainty-penalized cost of
Section~\ref{subsec:gp_uncertainty_cost}. Fifth, and directly motivated
by the inconclusive $\kappa$-sensitivity attempt above, diagnosing
\emph{why} the GP-uncertainty-penalized cost drives the SQP solver to
infeasibility independent of $\kappa$'s value — whether through
supplying an analytical (rather than finite-difference) Jacobian for
\texttt{cost\_gp.m} specifically, relaxing the control horizon
$N_c=1$ constraint identified as a contributing factor in the Stage~3
Jacobian test, or reformulating the uncertainty term as a tightened
constraint rather than a cost penalty (in the spirit of
\citet{hewing2020}'s probabilistic reachable sets) — would need to
precede any claim that $\kappa=0.8$ is a validated, robust operating
point rather than simply the value used throughout this paper's main
results. A structurally different route to the same real-time
feasibility goal, not explored in this paper, would replace the
nonlinear surrogate entirely with a Koopman-operator-based linear
approximation \citep{williams2015}, trading the accuracy of a
nonlinear model for the computational simplicity of linear MPC.
Coupling the corrected two-state
model to a higher-fidelity aeroelastic code (e.g.\ OpenFAST) is a longer-
term validation step once these nearer-term extensions are in place.
